\section{Scriptural Date and Location}
\zlabel{triptych:concise:chronology:start}
% Generated. Do not edit.
% Deterministic projection of research/chronology.toml.
% schema: 5
% document: liturgy/roman-rite/1962/propers/temporal/66-resumed-fifth-after-epiphany
% generated-by: tools/tpt proper-chronology annotations
\newcommand{\chronologyannotationclaim}[6]{#6}
\newcommand{\chronologyannotationcomparisonclaim}[8]{#8}
\newcommand{\chronologyannotationreach}[3]{}
\newcommand{\chronologyannotationgroup}[3]{#3}
\newcommand{\chronologyannotationcomparisongroup}[4]{#4}
\newcommand{\chronologyannotation}[1]{%
  \ifcsname triptychchronologyannotation@#1\endcsname
    \csname triptychchronologyannotation@#1\endcsname
  \else
    \PackageError{triptych}{No chronology annotation for #1}{}%
  \fi}
% comparison-dependency: src/gpt/liturgy/roman-rite/1962/propers/temporal/66-resumed-fifth-after-epiphany/research/chronology-profile-comparisons.toml sha256=0556d3276a9de8069099044e7830e0ad093924429fc6375de60c4481511e2788
% introit: Introit
\expandafter\def\csname triptychchronologyannotation@introit\endcsname{%
  \chronologyannotationgroup{historical-setting}{research-pending}{\textbf{Historical setting} -- Occasion date unresolved.}%
  \space \chronologyannotationgroup{composition}{nonuniform}{\textbf{Composition} -- No single date applies across every cited locus.}%
}
% epistle: Epistle
\expandafter\def\csname triptychchronologyannotation@epistle\endcsname{%
  \chronologyannotationgroup{composition}{disputed}{\textbf{Composition} -- disputed: \chronologyannotationreach{Col.3.12}{Col}{true}\chronologyannotationreach{Col.3.13}{Col}{true}\chronologyannotationreach{Col.3.14}{Col}{true}\chronologyannotationreach{Col.3.15}{Col}{true}\chronologyannotationreach{Col.3.16}{Col}{true}\chronologyannotationreach{Col.3.17}{Col}{true}\chronologyannotationclaim{composition.epistle-to-the-colossians}{composition}{catholic-traditional-v1}{disputed}{(Philemon; Colossians; Ephesians; Philippians), 61}{A.D. 61}; \chronologyannotationreach{Col.3.12}{Col}{true}\chronologyannotationreach{Col.3.13}{Col}{true}\chronologyannotationreach{Col.3.14}{Col}{true}\chronologyannotationreach{Col.3.15}{Col}{true}\chronologyannotationreach{Col.3.16}{Col}{true}\chronologyannotationreach{Col.3.17}{Col}{true}\chronologyannotationclaim{composition.epistle-to-the-colossians}{composition}{catholic-traditional-v1}{disputed}{some time between A.D. 61 and 63}{A.D. 61--63}.}%
}
% gradual: Gradual
\expandafter\def\csname triptychchronologyannotation@gradual\endcsname{%
  \chronologyannotationgroup{historical-setting}{research-pending}{\textbf{Historical setting} -- Occasion date unresolved.}%
  \space \chronologyannotationgroup{composition}{preferred}{\textbf{Composition}: \chronologyannotationreach{Ps.43.8}{Ps.41 Ps.42 Ps.43 Ps.44 Ps.45 Ps.46 Ps.47 Ps.48 Ps.83 Ps.84 Ps.86 Ps.87}{true}\chronologyannotationreach{Ps.43.9}{Ps.41 Ps.42 Ps.43 Ps.44 Ps.45 Ps.46 Ps.47 Ps.48 Ps.83 Ps.84 Ps.86 Ps.87}{true}\chronologyannotationclaim{composition.psalms-of-the-sons-of-korah}{composition}{catholic-traditional-v1}{preferred}{between the days of Isaias and the return from exile}{Between the days of Isaias and the return from exile}.}%
}
% alleluia: Alleluia
\expandafter\def\csname triptychchronologyannotation@alleluia\endcsname{%
  \chronologyannotationgroup{historical-setting}{research-pending}{\textbf{Historical setting} -- Occasion date unresolved.}%
  \space \chronologyannotationgroup{composition}{preferred}{\textbf{Composition}: \chronologyannotationreach{Ps.129.1}{Ps}{true}\chronologyannotationreach{Ps.129.2}{Ps}{true}\chronologyannotationclaim{critical.psalms.latest-composition-boundary}{composition}{catholic-critical-v1}{preferred}{before the Maccabean period, around 165 B.C.; no individual psalm can be dated securely}{Before c. 165 B.C.}}%
}
% gospel: Gospel
\expandafter\def\csname triptychchronologyannotation@gospel\endcsname{%
  \chronologyannotationgroup{narrated-event}{research-pending}{\textbf{Event} -- Narrated event date unresolved.}%
  \space \chronologyannotationgroup{composition}{disputed}{\textbf{Composition} -- disputed: \chronologyannotationreach{Matt.13.24}{Matt}{true}\chronologyannotationreach{Matt.13.25}{Matt}{true}\chronologyannotationreach{Matt.13.26}{Matt}{true}\chronologyannotationreach{Matt.13.27}{Matt}{true}\chronologyannotationreach{Matt.13.28}{Matt}{true}\chronologyannotationreach{Matt.13.29}{Matt}{true}\chronologyannotationreach{Matt.13.30}{Matt}{true}\chronologyannotationclaim{composition.gospel-of-matthew}{composition}{catholic-traditional-v1}{disputed}{about A.D. 38-45}{c. A.D. 38--45}; \chronologyannotationreach{Matt.13.24}{Matt}{true}\chronologyannotationreach{Matt.13.25}{Matt}{true}\chronologyannotationreach{Matt.13.26}{Matt}{true}\chronologyannotationreach{Matt.13.27}{Matt}{true}\chronologyannotationreach{Matt.13.28}{Matt}{true}\chronologyannotationreach{Matt.13.29}{Matt}{true}\chronologyannotationreach{Matt.13.30}{Matt}{true}\chronologyannotationclaim{composition.gospel-of-matthew}{composition}{catholic-traditional-v1}{disputed}{about the year 40-42}{c. A.D. 40--42}; \chronologyannotationreach{Matt.13.24}{Matt}{true}\chronologyannotationreach{Matt.13.25}{Matt}{true}\chronologyannotationreach{Matt.13.26}{Matt}{true}\chronologyannotationreach{Matt.13.27}{Matt}{true}\chronologyannotationreach{Matt.13.28}{Matt}{true}\chronologyannotationreach{Matt.13.29}{Matt}{true}\chronologyannotationreach{Matt.13.30}{Matt}{true}\chronologyannotationclaim{composition.gospel-of-matthew}{composition}{catholic-traditional-v1}{disputed}{the years 40-45}{A.D. 40--45}; \chronologyannotationreach{Matt.13.24}{Matt}{true}\chronologyannotationreach{Matt.13.25}{Matt}{true}\chronologyannotationreach{Matt.13.26}{Matt}{true}\chronologyannotationreach{Matt.13.27}{Matt}{true}\chronologyannotationreach{Matt.13.28}{Matt}{true}\chronologyannotationreach{Matt.13.29}{Matt}{true}\chronologyannotationreach{Matt.13.30}{Matt}{true}\chronologyannotationclaim{composition.gospel-of-matthew}{composition}{catholic-traditional-v1}{disputed}{about the year 60-68}{c. A.D. 60--68}; \chronologyannotationreach{Matt.13.24}{Matt}{true}\chronologyannotationreach{Matt.13.25}{Matt}{true}\chronologyannotationreach{Matt.13.26}{Matt}{true}\chronologyannotationreach{Matt.13.27}{Matt}{true}\chronologyannotationreach{Matt.13.28}{Matt}{true}\chronologyannotationreach{Matt.13.29}{Matt}{true}\chronologyannotationreach{Matt.13.30}{Matt}{true}\chronologyannotationclaim{composition.gospel-of-matthew}{composition}{catholic-traditional-v1}{disputed}{about the years 64-67}{c. A.D. 64--67}; \chronologyannotationreach{Matt.13.24}{Matt}{true}\chronologyannotationreach{Matt.13.25}{Matt}{true}\chronologyannotationreach{Matt.13.26}{Matt}{true}\chronologyannotationreach{Matt.13.27}{Matt}{true}\chronologyannotationreach{Matt.13.28}{Matt}{true}\chronologyannotationreach{Matt.13.29}{Matt}{true}\chronologyannotationreach{Matt.13.30}{Matt}{true}\chronologyannotationclaim{composition.gospel-of-matthew}{composition}{catholic-traditional-v1}{disputed}{about the year 50}{c. A.D. 50}.}%
  \space \chronologyannotationcomparisongroup{catholic-critical-v1}{composition}{composition-only}{\textbf{Composition (Catholic critical chronology, for comparison)}: \chronologyannotationreach{Matt.13.24}{Matt}{true}\chronologyannotationreach{Matt.13.25}{Matt}{true}\chronologyannotationreach{Matt.13.26}{Matt}{true}\chronologyannotationreach{Matt.13.27}{Matt}{true}\chronologyannotationreach{Matt.13.28}{Matt}{true}\chronologyannotationreach{Matt.13.29}{Matt}{true}\chronologyannotationreach{Matt.13.30}{Matt}{true}\chronologyannotationcomparisonclaim{critical.gospel-of-matthew}{composition}{catholic-critical-v1}{catholic-critical-v1}{preferred}{passage.united-states-conference-of-catholic-bishops.new-american-bible-revised-edition.english-usccb-web-2026-09-21.matthew-introduction}{post-A.D. 70 date}{The Greek Gospel of Matthew in the NABRE introduction, post-A.D. 70 date}.}%
}
% offertory: Offertory
\expandafter\def\csname triptychchronologyannotation@offertory\endcsname{%
  \chronologyannotationgroup{historical-setting}{research-pending}{\textbf{Historical setting} -- Occasion date unresolved.}%
  \space \chronologyannotationgroup{composition}{preferred}{\textbf{Composition}: \chronologyannotationreach{Ps.129.1}{Ps}{true}\chronologyannotationreach{Ps.129.2}{Ps}{true}\chronologyannotationclaim{critical.psalms.latest-composition-boundary}{composition}{catholic-critical-v1}{preferred}{before the Maccabean period, around 165 B.C.; no individual psalm can be dated securely}{Before c. 165 B.C.}}%
}
% communion: Communion
\expandafter\def\csname triptychchronologyannotation@communion\endcsname{%
  \chronologyannotationgroup{composition}{disputed}{\textbf{Composition} -- disputed: \chronologyannotationreach{Mark.11.24}{Mark}{true}\chronologyannotationclaim{composition.gospel-of-mark}{composition}{catholic-traditional-v1}{disputed}{the tenth or twelfth year after the Ascension (A.D. 38-40)}{A.D. 38--40}; \chronologyannotationreach{Mark.11.24}{Mark}{true}\chronologyannotationclaim{composition.gospel-of-mark}{composition}{catholic-traditional-v1}{disputed}{A.D. 40}{A.D. 40}; \chronologyannotationreach{Mark.11.24}{Mark}{true}\chronologyannotationclaim{composition.gospel-of-mark}{composition}{catholic-traditional-v1}{disputed}{the third year of Claudius (A.D. 43)}{A.D. 43}; \chronologyannotationreach{Mark.11.24}{Mark}{true}\chronologyannotationclaim{composition.gospel-of-mark}{composition}{catholic-traditional-v1}{disputed}{somewhere between A.D. 50 and 67}{A.D. 50--67}; \chronologyannotationreach{Mark.11.24}{Mark}{true}\chronologyannotationclaim{composition.gospel-of-mark}{composition}{catholic-traditional-v1}{disputed}{between 60 and 70}{A.D. 60--70}; \chronologyannotationreach{Mark.11.24}{Mark}{true}\chronologyannotationclaim{composition.gospel-of-mark}{composition}{catholic-traditional-v1}{disputed}{after Matthew, and not as late as the destruction of Jerusalem}{After Matthew, and not as late as the destruction of Jerusalem}.}%
  \space \chronologyannotationcomparisongroup{catholic-critical-v1}{composition}{composition-only}{\textbf{Composition (Catholic critical chronology, for comparison)}: \chronologyannotationreach{Mark.11.24}{Mark.1 Mark.2 Mark.3 Mark.4 Mark.5 Mark.6 Mark.7 Mark.8 Mark.9 Mark.10 Mark.11 Mark.12 Mark.13 Mark.14 Mark.15 Mark.16.1-8}{true}\chronologyannotationcomparisonclaim{critical.gospel-of-mark}{composition}{catholic-critical-v1}{catholic-critical-v1}{preferred}{artifact.united-states-conference-of-catholic-bishops.new-american-bible-revised-edition.english-usccb-web-2026-11-02.mark-introduction-web-text-272bcce7}{perhaps shortly after the year 70}{The Gospel of Mark under the NABRE's tentative modern dating proposal, perhaps shortly after A.D. 70}.}%
}
% presentation: orientation-v2
\newcommand{\chronologyorientationrow}[2]{%
  \ifcsname triptychchronologyorientation@#1\endcsname
    \csname triptychchronologyorientation@#1\endcsname{#2}%
  \else
    \PackageError{triptych}{No chronology orientation row for #1}{}%
  \fi}
\expandafter\def\csname triptychchronologyorientation@row-4497fc13bd16e15c\endcsname#1{%
Introit: Ps 84:2 (Vulgate numbering; Hebrew 85); Gradual: Psalm 43:8--9 (Vulgate numbering; Hebrew 44); Alleluia: Psalm 129:1--2 (Vulgate numbering; Hebrew 130); Offertory: Psalm 129:1--2 (Vulgate numbering; Hebrew 130) & Historical setting. #1 & Occasion date unresolved\\[.25em]%
}
\expandafter\def\csname triptychchronologyorientation@row-d1f3edd1ff63c2d2\endcsname#1{%
Introit: Ps 84:2 (Vulgate numbering; Hebrew 85); Gradual: Psalm 43:8--9 (Vulgate numbering; Hebrew 44) & Composition: The Psalms of the Sons of Korah. #1 & \chronologyannotationreach{Ps.43.8}{Ps.41 Ps.42 Ps.43 Ps.44 Ps.45 Ps.46 Ps.47 Ps.48 Ps.83 Ps.84 Ps.86 Ps.87}{true}\chronologyannotationreach{Ps.43.9}{Ps.41 Ps.42 Ps.43 Ps.44 Ps.45 Ps.46 Ps.47 Ps.48 Ps.83 Ps.84 Ps.86 Ps.87}{true}\chronologyannotationreach{Ps.84.2}{Ps.41 Ps.42 Ps.43 Ps.44 Ps.45 Ps.46 Ps.47 Ps.48 Ps.83 Ps.84 Ps.86 Ps.87}{true}\chronologyannotationclaim{composition.psalms-of-the-sons-of-korah}{composition}{catholic-traditional-v1}{preferred}{between the days of Isaias and the return from exile}{Between the days of Isaias and the return from exile}\\[.25em]%
}
\expandafter\def\csname triptychchronologyorientation@row-da4404f3dbb16b26\endcsname#1{%
Alleluia: Psalm 129:1--2 (Vulgate numbering; Hebrew 130); Offertory: Psalm 129:1--2 (Vulgate numbering; Hebrew 130) & Composition: Psalms. #1 & \chronologyannotationreach{Ps.129.1}{Ps}{true}\chronologyannotationreach{Ps.129.2}{Ps}{true}\chronologyannotationclaim{critical.psalms.latest-composition-boundary}{composition}{catholic-critical-v1}{preferred}{before the Maccabean period, around 165 B.C.; no individual psalm can be dated securely}{Before c. 165 B.C.}\\[.25em]%
}
\expandafter\def\csname triptychchronologyorientation@row-84a964a3effacafd\endcsname#1{%
Introit: Jer 29:11--12,14 & Composition: Jeremias's letter to the exiles after Jechonias's deportation. #1 & \chronologyannotationreach{Jer.29.11}{Jer.29.1-23}{false}\chronologyannotationreach{Jer.29.12}{Jer.29.1-23}{false}\chronologyannotationreach{Jer.29.14}{Jer.29.1-23}{false}\chronologyannotationclaim{composition.jeremias-letter-to-the-exiles}{composition}{catholic-traditional-v1}{preferred}{After that Jechonias the king ... were departed out of Jerusalem}{After Jechonias and the deportees had departed from Jerusalem}\\[.25em]%
}
\expandafter\def\csname triptychchronologyorientation@row-cc5afc871115711e\endcsname#1{%
Introit: Jer 29:11--12,14 & Final formation: The received Book of Jeremias after enlargement. #1 & \chronologyannotationreach{Jer.29.11}{Jer}{true}\chronologyannotationreach{Jer.29.12}{Jer}{true}\chronologyannotationreach{Jer.29.14}{Jer}{true}\chronologyannotationclaim{composition.book-of-jeremias}{final-formation}{catholic-traditional-v1}{preferred}{the first transcript was enlarged by various insertions and additions}{After the first authentic recension}\\[.25em]%
}
\expandafter\def\csname triptychchronologyorientation@row-a1420e3e43d2b570\endcsname#1{%
Gospel: Matthew 13:24--30 & Event. #1 & Date unresolved\\[.25em]%
}
\expandafter\def\csname triptychchronologyorientation@row-68d3fd83119cd2be\endcsname#1{%
Gospel: Matthew 13:24--30 & Composition: The Gospel of St Matthew. #1 & Disputed: \chronologyannotationreach{Matt.13.24}{Matt}{true}\chronologyannotationreach{Matt.13.25}{Matt}{true}\chronologyannotationreach{Matt.13.26}{Matt}{true}\chronologyannotationreach{Matt.13.27}{Matt}{true}\chronologyannotationreach{Matt.13.28}{Matt}{true}\chronologyannotationreach{Matt.13.29}{Matt}{true}\chronologyannotationreach{Matt.13.30}{Matt}{true}\chronologyannotationclaim{composition.gospel-of-matthew}{composition}{catholic-traditional-v1}{disputed}{about A.D. 38-45}{c. A.D. 38--45}; \chronologyannotationreach{Matt.13.24}{Matt}{true}\chronologyannotationreach{Matt.13.25}{Matt}{true}\chronologyannotationreach{Matt.13.26}{Matt}{true}\chronologyannotationreach{Matt.13.27}{Matt}{true}\chronologyannotationreach{Matt.13.28}{Matt}{true}\chronologyannotationreach{Matt.13.29}{Matt}{true}\chronologyannotationreach{Matt.13.30}{Matt}{true}\chronologyannotationclaim{composition.gospel-of-matthew}{composition}{catholic-traditional-v1}{disputed}{about the year 40-42}{c. A.D. 40--42}; \chronologyannotationreach{Matt.13.24}{Matt}{true}\chronologyannotationreach{Matt.13.25}{Matt}{true}\chronologyannotationreach{Matt.13.26}{Matt}{true}\chronologyannotationreach{Matt.13.27}{Matt}{true}\chronologyannotationreach{Matt.13.28}{Matt}{true}\chronologyannotationreach{Matt.13.29}{Matt}{true}\chronologyannotationreach{Matt.13.30}{Matt}{true}\chronologyannotationclaim{composition.gospel-of-matthew}{composition}{catholic-traditional-v1}{disputed}{the years 40-45}{A.D. 40--45}; \chronologyannotationreach{Matt.13.24}{Matt}{true}\chronologyannotationreach{Matt.13.25}{Matt}{true}\chronologyannotationreach{Matt.13.26}{Matt}{true}\chronologyannotationreach{Matt.13.27}{Matt}{true}\chronologyannotationreach{Matt.13.28}{Matt}{true}\chronologyannotationreach{Matt.13.29}{Matt}{true}\chronologyannotationreach{Matt.13.30}{Matt}{true}\chronologyannotationclaim{composition.gospel-of-matthew}{composition}{catholic-traditional-v1}{disputed}{about the year 60-68}{c. A.D. 60--68}; \chronologyannotationreach{Matt.13.24}{Matt}{true}\chronologyannotationreach{Matt.13.25}{Matt}{true}\chronologyannotationreach{Matt.13.26}{Matt}{true}\chronologyannotationreach{Matt.13.27}{Matt}{true}\chronologyannotationreach{Matt.13.28}{Matt}{true}\chronologyannotationreach{Matt.13.29}{Matt}{true}\chronologyannotationreach{Matt.13.30}{Matt}{true}\chronologyannotationclaim{composition.gospel-of-matthew}{composition}{catholic-traditional-v1}{disputed}{about the years 64-67}{c. A.D. 64--67}; \chronologyannotationreach{Matt.13.24}{Matt}{true}\chronologyannotationreach{Matt.13.25}{Matt}{true}\chronologyannotationreach{Matt.13.26}{Matt}{true}\chronologyannotationreach{Matt.13.27}{Matt}{true}\chronologyannotationreach{Matt.13.28}{Matt}{true}\chronologyannotationreach{Matt.13.29}{Matt}{true}\chronologyannotationreach{Matt.13.30}{Matt}{true}\chronologyannotationclaim{composition.gospel-of-matthew}{composition}{catholic-traditional-v1}{disputed}{about the year 50}{c. A.D. 50}\\[.25em]%
}
\expandafter\def\csname triptychchronologyorientation@row-941543664efaa9ec\endcsname#1{%
Gospel: Matthew 13:24--30 & Composition (Catholic critical chronology, for comparison): The Greek Gospel of Matthew in the NABRE introduction. #1 & \chronologyannotationreach{Matt.13.24}{Matt}{true}\chronologyannotationreach{Matt.13.25}{Matt}{true}\chronologyannotationreach{Matt.13.26}{Matt}{true}\chronologyannotationreach{Matt.13.27}{Matt}{true}\chronologyannotationreach{Matt.13.28}{Matt}{true}\chronologyannotationreach{Matt.13.29}{Matt}{true}\chronologyannotationreach{Matt.13.30}{Matt}{true}\chronologyannotationcomparisonclaim{critical.gospel-of-matthew}{composition}{catholic-critical-v1}{catholic-critical-v1}{preferred}{passage.united-states-conference-of-catholic-bishops.new-american-bible-revised-edition.english-usccb-web-2026-09-21.matthew-introduction}{post-A.D. 70 date}{Post-A.D. 70 date}\\[.25em]%
}
\expandafter\def\csname triptychchronologyorientation@row-8b7f20d842814860\endcsname#1{%
Communion: Mark 11:24 & Composition: The Gospel of St Mark. #1 & Disputed: \chronologyannotationreach{Mark.11.24}{Mark}{true}\chronologyannotationclaim{composition.gospel-of-mark}{composition}{catholic-traditional-v1}{disputed}{the tenth or twelfth year after the Ascension (A.D. 38-40)}{A.D. 38--40}; \chronologyannotationreach{Mark.11.24}{Mark}{true}\chronologyannotationclaim{composition.gospel-of-mark}{composition}{catholic-traditional-v1}{disputed}{A.D. 40}{A.D. 40}; \chronologyannotationreach{Mark.11.24}{Mark}{true}\chronologyannotationclaim{composition.gospel-of-mark}{composition}{catholic-traditional-v1}{disputed}{the third year of Claudius (A.D. 43)}{A.D. 43}; \chronologyannotationreach{Mark.11.24}{Mark}{true}\chronologyannotationclaim{composition.gospel-of-mark}{composition}{catholic-traditional-v1}{disputed}{somewhere between A.D. 50 and 67}{A.D. 50--67}; \chronologyannotationreach{Mark.11.24}{Mark}{true}\chronologyannotationclaim{composition.gospel-of-mark}{composition}{catholic-traditional-v1}{disputed}{between 60 and 70}{A.D. 60--70}; \chronologyannotationreach{Mark.11.24}{Mark}{true}\chronologyannotationclaim{composition.gospel-of-mark}{composition}{catholic-traditional-v1}{disputed}{after Matthew, and not as late as the destruction of Jerusalem}{After Matthew, and not as late as the destruction of Jerusalem}\\[.25em]%
}
\expandafter\def\csname triptychchronologyorientation@row-f35007a3f1f636d7\endcsname#1{%
Communion: Mark 11:24 & Composition (Catholic critical chronology, for comparison): The Gospel of Mark under the NABRE's tentative modern dating proposal. #1 & \chronologyannotationreach{Mark.11.24}{Mark.1 Mark.2 Mark.3 Mark.4 Mark.5 Mark.6 Mark.7 Mark.8 Mark.9 Mark.10 Mark.11 Mark.12 Mark.13 Mark.14 Mark.15 Mark.16.1-8}{true}\chronologyannotationcomparisonclaim{critical.gospel-of-mark}{composition}{catholic-critical-v1}{catholic-critical-v1}{preferred}{artifact.united-states-conference-of-catholic-bishops.new-american-bible-revised-edition.english-usccb-web-2026-11-02.mark-introduction-web-text-272bcce7}{perhaps shortly after the year 70}{Perhaps shortly after A.D. 70}\\[.25em]%
}
\expandafter\def\csname triptychchronologyorientation@row-6be86a5d4d4ef64b\endcsname#1{%
Epistle: Colossians 3:12--17 & Composition: The Epistle to the Colossians. #1 & Disputed: \chronologyannotationreach{Col.3.12}{Col}{true}\chronologyannotationreach{Col.3.13}{Col}{true}\chronologyannotationreach{Col.3.14}{Col}{true}\chronologyannotationreach{Col.3.15}{Col}{true}\chronologyannotationreach{Col.3.16}{Col}{true}\chronologyannotationreach{Col.3.17}{Col}{true}\chronologyannotationclaim{composition.epistle-to-the-colossians}{composition}{catholic-traditional-v1}{disputed}{(Philemon; Colossians; Ephesians; Philippians), 61}{A.D. 61}; \chronologyannotationreach{Col.3.12}{Col}{true}\chronologyannotationreach{Col.3.13}{Col}{true}\chronologyannotationreach{Col.3.14}{Col}{true}\chronologyannotationreach{Col.3.15}{Col}{true}\chronologyannotationreach{Col.3.16}{Col}{true}\chronologyannotationreach{Col.3.17}{Col}{true}\chronologyannotationclaim{composition.epistle-to-the-colossians}{composition}{catholic-traditional-v1}{disputed}{some time between A.D. 61 and 63}{A.D. 61--63}\\[.25em]%
}

\begin{chronologyorientationtable}
\chronologyorientationrow{row-4497fc13bd16e15c}{Israel recalls rescue and seeks mercy; exact occasions and places unknown.}
\chronologyorientationrow{row-d1f3edd1ff63c2d2}{Temple singers; Jerusalem worship horizon, writing sites unknown.}
\chronologyorientationrow{row-da4404f3dbb16b26}{Israel hopes for redemption; writing site unknown. Bound is Psalter-wide.}
\chronologyorientationrow{row-84a964a3effacafd}{Jerusalem to Babylon; Jer 29:1--3.}
\chronologyorientationrow{row-cc5afc871115711e}{Judah/Jerusalem ministry; final-redaction site unknown.}
\chronologyorientationrow{row-a1420e3e43d2b570}{Crowds beside the sea; Matt 13:1--3.}
\chronologyorientationrow{row-68d3fd83119cd2be}{Traditional original associated with departure from Palestine.}
\chronologyorientationrow{row-941543664efaa9ec}{Probable majority proposal; Antioch plausible.}
\chronologyorientationrow{row-8b7f20d842814860}{Traditional Rome; instruction in Jerusalem/Bethania narrative.}
\chronologyorientationrow{row-f35007a3f1f636d7}{Tentative anonymous-author proposal; Syria possible.}
\chronologyorientationrow{row-6be86a5d4d4ef64b}{Traditional Roman captivity; recipients at Colossae. Modern account leaves authorship and prison site uncertain.}
\end{chronologyorientationtable}
{\footnotesize Full source loci and distinctions: expansive study, Historical appendix, ``Psalms,'' ``Jeremias,'' ``Matthew,'' ``Mark'' and ``Colossians.'' Gospel composition dates concern writing; the teaching scene is a separate subject.}
\zlabel{triptych:concise:chronology:end}
